\section{格点涂抹矩阵元的重正化}\label{sec:smearing}

在许多格点计算中，数据的噪声可能很大，因而需要更多的统计量。然而大统计意味着更多资源，有时难以获得。为此，格点工作者发明了压制短距离涨落的唯象手段：涂抹。但涂抹可能导致所处理的组态与基本理论脱节。另一方面，从某种意义上说，涂抹只是让有效紫外截断（格点间距的倒数）变小；而在 $a\to 0$ 极限下，一切涂抹都成为定域效应，可以通过某种重正化恰当地计入。本节检验前一节讨论的流程对涂抹矩阵元是否仍然有效。

\begin{figure*}[tbp]
\centering
\includegraphics[width=8.8cm]{images/fitting_overlap_quark_hyp1_Oa.pdf}
\includegraphics[width=8.8cm]{images/fitting_clover_quark_hyp1_Oa.pdf}
\includegraphics[width=8.8cm]{images/fitting_overlap_Pion_hyp1_Oa.pdf}
\includegraphics[width=8.8cm]{images/fitting_clover_Pion_hyp1_Oa.pdf}
\includegraphics[width=8.8cm]{images/fitting_clover_Nucleon_hyp1_Oa.pdf}
\caption{用式~(\ref{eq:logM_ss}) 拟合 HYP 涂抹情形的 $O_{\gamma_t}(z)$ 矩阵元。$\Lambda_{\rm QCD}$ 固定为 0.39 GeV。
}
\label{fig:fitting_hyp}
\end{figure*}


\begin{figure*}[tbp]
\centering
\includegraphics[width=8.5cm]{images/ez_vs_z_overlap_quark_hyp1_Oa.pdf}
\includegraphics[width=8.5cm]{images/ez_vs_z_clover_quark_hyp1_Oa.pdf}
\includegraphics[width=8.5cm]{images/ez_vs_z_overlap_Pion_hyp1_Oa.pdf}
\includegraphics[width=8.5cm]{images/ez_vs_z_clover_Pion_hyp1_Oa.pdf}
\includegraphics[width=8.5cm]{images/ez_vs_z_clover_Nucleon_hyp1_Oa.pdf}
\caption{$e(z)$ 随 $z$ 的变化。$e(z)$ 由图~\ref{fig:fitting_hyp} 中的拟合提取，斜率标注于图例。
}
\label{fig:e(z)_hyp}
\end{figure*}

为确保 HYP 涂抹数据可用于精细的分析，我们可以像在第~\ref{sec:test} 节中那样做一个简单的检验：考察 HYP 涂抹数据是否呈现微扰论预言的线性发散性质。这里用如下函数拟合裸矩阵元 $\mathcal{M}$ 以提取线性发散因子，
\begin{align}\label{eq:logM_ss}
\ln \mathcal{M}(z,a) = \frac{e(z)}{a \ln[a \Lambda_{\rm QCD}]}
+ g(z) + f_{1,2}(z)a,
\end{align}
其中允许存在离散化误差以改善拟合。在这个简单测试中我们把 $\Lambda_{\rm QCD}$ 固定为 0.39 GeV。图~\ref{fig:fitting_hyp} 展示了不同作用量与不同态的拟合，它们看上去都是合理的。
$\chi^2$ 未在此展示，将在提取重正化因子时予以考虑。

图~\ref{fig:e(z)_hyp} 展示了提取出的 $e(z)$ 随 $z$ 的变化。线性的 $z$ 依赖大体重现，但在大 $z$ 处可见一些偏离。我们可以对 $e(z)$ 作线性拟合并比较拟合斜率。
总体而言，拟合质量不如未涂抹情形，但仍相当可观。
从 clover 夸克提取到的线性发散斜率与其他情形完全不同，这解释了为什么对 HYP 涂抹数据、clover 情形不能用 RI/MOM 因子消除线性发散~\cite{Zhang:2020rsx}。虽然 overlap 夸克与 $\pi$、clover $\pi$ 与核子的斜率彼此相近，但仍存在一些小出入。对这些情形使用 RI/MOM 因子可能残留小的线性发散，而自重正化方法为矩阵元重正化提供了更好的途径。

接下来，我们把自重正化方法推广到 HYP 涂抹数据。我们的拟合函数式~(\ref{eq:logM}) 对 HYP 涂抹数据未必有明确的物理含义，
但可以把它看作一个唯象模型。

\begin{figure*}[tbp]
\centering
\includegraphics[width=8.5cm]{images/chisqmap_overlap_quark_hyp1.pdf}
\includegraphics[width=8.5cm]{images/chisqmap_clover_quark_hyp1.pdf}
\includegraphics[width=8.5cm]{images/chisqmap_overlap_Pion_hyp1.pdf}
\includegraphics[width=8.5cm]{images/chisqmap_clover_Pion_hyp1.pdf}
\includegraphics[width=8.5cm]{images/chisqmap_clover_Nucleon_hyp1.pdf}
\caption{用式~(\ref{eq:logM}) 拟合 HYP 涂抹的 $O_{\gamma_t}(z)$ 矩阵元。图中给出不同作用量与不同态下 $\chi^{2}$ 随 $k$ 与 $\Lambda_{\rm QCD}$ 的分布。$\langle  \chi^{2}/{\rm d.o.f.} \rangle_{z}$ 是各 $z$ 拟合的 $\chi^{2}/{\rm d.o.f.}$ 的平均值。
}
\label{fig:chisqmap_hyp}
\end{figure*}


\begin{figure*}[tbp]
\centering
\includegraphics[width=8.5cm]{images/Comp_overlap_quark.pdf}
\includegraphics[width=8.5cm]{images/Comp_clover_quark.pdf}
\includegraphics[width=8.5cm]{images/Comp_overlap_Pion.pdf}
\includegraphics[width=8.5cm]{images/Comp_clover_Pion.pdf}
\includegraphics[width=8.5cm]{images/Comp_clover_Nucleon.pdf}
\caption{未涂抹（hyp0，红）与 HYP 涂抹（hyp1，蓝）两种情形下重正化矩阵元 Exp[$g(z)-m_{0}z$] 的比较。
}
\label{fig:Re_hyp}
\end{figure*}

我们先在线性发散部分采用带两个参数（$k$ 与 $\Lambda_{\rm QCD}$）的拟合函数。
图~\ref{fig:chisqmap_hyp} 给出了 $\chi^{2}$ 随它们的分布。由图可以清楚地看到，在“小卡方带”内再也找不到 $k$ 取单圈微扰值 7.4 GeV$^{-1}$fm$^{-1}$（无量纲时为 1.46，见式~(\ref{eq:k})）的解。
这在预期之内：格点“空间”被放大约一倍，因此可以预期“小卡方带”内 $k$ 的取值至多约为其一半，即 3.7 GeV$^{-1}$fm$^{-1}$（无量纲时为 0.73）。实际上，最概然的 $k$ 现在大约在 3.0 GeV$^{-1}$fm$^{-1}$（无量纲时为 0.59）或以下。换一种说法：HYP 涂抹把线性发散这类短程行为平滑掉了。另一方面，$\Lambda_{\rm QCD}$ 的取值范围更大，可能再次对应更大的有效格点间距。

于是对 HYP 涂抹情形，我们在“小卡方带”中心附近选取一组 ($k$, $\Lambda_{\rm QCD}$)。对 overlap 夸克与 clover 夸克取 $k=2.5$ GeV$^{-1}$fm$^{-1}$（无量纲时为 0.49），对 overlap $\pi$、clover $\pi$ 和 clover 核子取 $k=2.6$ GeV$^{-1}$fm$^{-1}$（无量纲时为 0.51）。相应的 $\Lambda_{\rm QCD}$ 分散较大：overlap 夸克与 $\pi$ 关联分别取 $\Lambda_{\rm QCD} = 0.43$ 与 $0.38$ GeV，clover 夸克、$\pi$ 与核子关联分别取 $\Lambda_{\rm QCD} = 0.61$、$0.37$ 与 $0.35$ GeV。

选定 ($k$, $\Lambda_{\rm QCD}$) 参数之后，我们减除线性发散，再把结果匹配到微扰矩阵元，以确定针对非微扰质量效应可能需要的进一步减除。图~\ref{fig:Re_hyp} 用蓝色标记和线条展示 HYP 涂抹下的重正化矩阵元，并与不经 HYP 涂抹的情形进行对比。

涂抹与非涂抹矩阵元之间的巨大差别出现在大 $z$ 处的 clover 夸克关联情形。
我们再次推测，由于手征对称性破缺效应，那里线性发散表现出某种奇特的行为。不过在 $\pi$ 矩阵元中差别较小，尽管其误差棒也更大。对 overlap 情形，夸克与 $\pi$ 的关联没有因涂抹效应而发生实质变化。另一方面，核子情形的矩阵元在涂抹后误差棒明显更小，这显示了涂抹的功效。

总之，我们用于重正化线性发散矩阵元的分析方法同样能成功地应用于涂抹矩阵元。涂抹似乎并未定性地改变重正化的行为，而只是定量地有所改变。
在（适度）涂抹下内禀物理似乎不变，这与“涂抹是降低格点计算统计噪声的一种替代手段”的预期相一致。
